\section{The chessboard estimate}

\begin{definition}[Reflection positivity, Georgii (17.6)--(17.7)]
    \label{def:rp-17-7}
    \uses{def:trans-transformation}
    \lean{MeasureTheory.GibbsMeasure.genReflection, MeasureTheory.GibbsMeasure.IsReflectionPositive}
    \leanok

    On the torus $\mathbb Z_{2N}$ the generalised reflection $\tilde r$ is the site reflection
    $z \mapsto -1-z$ composed with a pure-spin involution. A measure $\mu$ is reflection positive
    if $\int f\,(f\circ\tilde r)\,d\mu \geq 0$ for every bounded $f$ measurable with respect to
    the positive half.
\end{definition}

\begin{lemma}[Georgii (17.8)]
    \label{lem:rp-17-8}
    \uses{def:rp-17-7}
    \lean{MeasureTheory.GibbsMeasure.sq_integral_mul_comp_le, sq_le_mul_of_forall_quadratic_nonneg}
    \leanok

    If $\mu$ is reflection positive and invariant under the involution $\tilde r$, then
    $\bigl(\int f\,(g\circ\tilde r)\,d\mu\bigr)^2 \leq \int f\,(f\circ\tilde r)\,d\mu
    \int g\,(g\circ\tilde r)\,d\mu$. Without invariance the statement is false: a nonnegative
    real form only bounds the symmetrised pairing.
\end{lemma}
\begin{proof}
    \leanok
    Expand $\int (f+tg)((f+tg)\circ\tilde r)\,d\mu \geq 0$ in $t$ and use the discriminant.
\end{proof}

\begin{lemma}[Georgii (17.9)]
    \label{lem:rp-17-9}
    \uses{lem:rp-17-8}
    \lean{MeasureTheory.GibbsMeasure.altConfig, MeasureTheory.GibbsMeasure.foldPos, MeasureTheory.GibbsMeasure.foldNeg, MeasureTheory.GibbsMeasure.altConfig_of_shift_of_foldPos, MeasureTheory.GibbsMeasure.prod_foldPos_mul_prod_foldNeg, MeasureTheory.GibbsMeasure.pow_le_prod_of_chessboard}
    \leanok

    Let $D$ be a nonnegative function of words of length $2N$ over a finite alphabet with an
    involution $t$, invariant under cyclic shift and satisfying
    $D(\alpha)^2 \leq D(\alpha^+)D(\alpha^-)$ for the two folded words. Then
    $D(\alpha)^{2N} \leq \prod_{z} D(\alpha_z, t\alpha_z, \alpha_z, \ldots)$.
\end{lemma}
\begin{proof}
    \leanok
    Georgii's two passes over the strips $R_{a,\ell}$, with the doubling step
    $\ell \mapsto N \wedge 2\ell$, are one inductive argument run twice: once to see that
    $D > 0$ propagates, once at a word maximising $D$ among words with the same letters. The
    normalising weights obey the exact identity
    $\prod Q(\alpha^+)\prod Q(\alpha^-) = (\prod Q(\alpha))^2$.
\end{proof}

\begin{theorem}[Georgii (17.11)]
    \label{thm:rp-17-11}
    \uses{lem:rp-17-9, def:rp-17-7}
    \lean{MeasureTheory.GibbsMeasure.wordProd, MeasureTheory.GibbsMeasure.sq_integral_wordProd_le, MeasureTheory.GibbsMeasure.integral_wordProd_shift, MeasureTheory.GibbsMeasure.abs_integral_prod_pow_le, MeasureTheory.GibbsMeasure.torusReflAt, MeasureTheory.GibbsMeasure.torusPosAt, MeasureTheory.GibbsMeasure.genReflectionAt, MeasureTheory.GibbsMeasure.tauPow, MeasureTheory.GibbsMeasure.IsReflectionPositive.map, MeasureTheory.GibbsMeasure.splitHead, MeasureTheory.GibbsMeasure.splitTail, MeasureTheory.GibbsMeasure.isReflectionPositive_map_splitHead, MeasureTheory.GibbsMeasure.isReflectionPositive_map_splitTail, MeasureTheory.GibbsMeasure.abs_integral_prod_pow_le_pi_abs, MeasureTheory.GibbsMeasure.integral_prod_tauPow_nonneg, MeasureTheory.GibbsMeasure.abs_integral_prod_pow_le_pi}
    \leanok

    Let $\Lambda = \mathbb Z_{2N}^d$, $d \geq 1$, and let $\mu$ be a finite measure on $E^\Lambda$
    which is $\Lambda$-periodic and, for every direction $k$, $\tilde r_k$-positive and
    $\tilde r_k$-invariant. For bounded measurable $f_i$ on $E$,
    $\bigl|\int \prod_i f_i(\omega_i)\,d\mu\bigr|^{|\Lambda|} \leq
    \prod_j \int \prod_i f_j(\tau^i\omega_i)\,d\mu$, the right side being nonnegative.
\end{theorem}
\begin{proof}
    \leanok
    In one dimension Lemma \ref{lem:rp-17-8} is hypothesis (ii) of Lemma \ref{lem:rp-17-9} for
    $D(\alpha) = |\int \prod_i \alpha_i(\omega_i)\,d\mu|$, and shift invariance is (i). In
    dimension $d + 1$ the torus is $\mathbb Z_{2N} \times \mathbb Z_{2N}^d$; viewing $\mu$ on
    $(E^{\mathbb Z_{2N}^d})^{\mathbb Z_{2N}}$ (the head coordinate) reduces to $d = 1$ with state
    space $E^{\mathbb Z_{2N}^d}$ and the coordinatewise involution, and viewing it on
    $(E^{\mathbb Z_{2N}})^{\mathbb Z_{2N}^d}$ (the tail) applies the induction hypothesis to the
    factors, since periodicity, positivity and invariance pass to both images.
\end{proof}
